\chapter{How pystencils 1.4 Vectorizes}\label{r04:ch:vectorizer}

This chapter describes the part of \ps{}~1.4 that every vector instruction set goes
through, RVV included. Its design explains both what upstream's \isa{rvv} looks like
(\cref{r04:ch:upstream-rvv}) and why our extras could be written without touching any
of it (\cref{r04:ch:rvv-ext}). \Cref{r04:fig:vectorizer} is the map. The files are in
\file{src/pystencils/} of the pystencils submodule at \code{release/1.4}.

\begin{figure}[htbp]
  \centering
  \resizebox{0.92\linewidth}{!}{% pystencils 1.4: from assignments to vectorized C code (chapter 4).
\begin{tikzpicture}[node distance=4mm and 7mm]
  \node[r04box,text width=40mm] (asm) {SymPy assignments\\[1pt]\scriptsize on abstract fields};
  \node[r04box,below=of asm,text width=40mm] (ast) {AST\\[1pt]\scriptsize loop nest, types,\\
    base pointers, constants\\hoisted before the loops};
  \node[r04box,below=of ast,text width=40mm] (omp) {blocking, OpenMP\\[1pt]\scriptsize (optional)};
  \node[r04box,below=of omp,text width=40mm,fill=codebg] (vec) {\code{vectorize()}\\[1pt]
    \scriptsize innermost loop: cut or keep,\\vector accesses, vector casts};
  \node[r04box,below=of vec,text width=40mm,fill=codebg] (print) {C printer\\[1pt]
    \scriptsize\code{VectorizedCustomSympyPrinter}};
  \node[r04box,below=of print,text width=40mm] (c) {C function};
  \foreach \a/\b in {asm/ast,ast/omp,omp/vec,vec/print,print/c} \draw[r04arrow] (\a) -- (\b);

  % the instruction set dictionary
  \node[r04code,right=28mm of vec,text width=64mm,align=left,yshift=-6mm] (dict) {%
    'width': vsetvlmax\_e32m1()\\
    '+': \_\_riscv\_vfadd\_vv\_f32m1(\{0\}, \{1\}, vl)\\
    'loadU': \_\_riscv\_vle32\_v\_f32m1(\{0\}, vl)\\
    'makeVecConst': \_\_riscv\_vfmv\_v\_f\_f32m1(\{0\}, vl)\\
    'cachelineZero': cachelineZero((void*) \{0\})\\
    \dots\\[2pt]
    \textrm{\scriptsize with } vl = \{loop\_stop\} - \{loop\_counter\}};
  \node[r04label,above=1mm of dict] (dictlabel) {instruction set = dictionary of C templates\\
    \code{get_vector_instruction_set(dtype, name)}};
  \draw[r04arrow,dashed] (dict.west) -- ++(-9mm,0) |- node[r04label,pos=0.78,above]{width,\\loop rule} (vec.east);
  \draw[r04arrow,dashed] (dict.west) -- ++(-9mm,0) |- node[r04label,pos=0.78,above]{\code{.format}} (print.east);

  \node[r04box,r04ours,below=8mm of dict,text width=64mm] (ext) {\isa{rvv} or \isa{rvv_ext-<options>}:\\
    \scriptsize the only part that differs between the two (chapters 5 and 6)};
  \draw[r04arrow] (ext) -- (dict);
\end{tikzpicture}
}
  \caption[pystencils 1.4: from assignments to vectorized C code]{\ps{}~1.4 from
    assignments to vectorized C code. All instruction sets share the pipeline on the
    left; an instruction set is a dictionary of C string templates (right) that the
    vectorizer reads for the vector width and the loop rule, and that the printer fills
    in. \texttt{rvv} and \texttt{rvv\_ext} differ only in that dictionary.}
  \label{r04:fig:vectorizer}
\end{figure}

\section{From assignments to an AST}

\code{create_kernel} (\file{kernelcreation.py}) checks the assignments and builds the
CPU kernel (\file{cpu/kernelcreation.py}): it types all symbols, wraps the assignments
in a loop nest over the fields' domain, resolves field accesses to pointer arithmetic
with one base pointer per field and row (the $y$, $z$ and index offsets are added
outside the innermost loop), and moves constant subexpressions in front of the loops.
Loop blocking and OpenMP follow if requested. Only then, on this finished scalar AST, it
calls the vectorizer:
\begin{pycode}
if config.cpu_vectorize_info:                                  # kernelcreation.py
    if isinstance(config.cpu_vectorize_info, dict):
        vectorize(ast, **config.cpu_vectorize_info)
\end{pycode}
Note what that means for constants: a hoisted constant is a \emph{scalar} declaration
before the loops (\code{const float xi_20 = 2.0f;}). Inside the vector loop it is used
in vector expressions, so the vectorizer wraps every use in a cast to a vector type,
and the printer turns each cast into a broadcast (\cref{r04:sec:printer}).

\section{The options}

\code{vectorize} (\file{cpu/vectorization.py}) takes five keyword options, the same
keys that waLBerla's \code{cpu_vectorize_info} carries (\cref{r04:sec:where-isa}):
\begin{description}[style=nextline]
  \item[\code{instruction_set}] a name such as \code{'avx'}, \code{'neon'},
    \code{'sve256'}, \code{'rvv'}; \code{'best'} takes the last entry of
    \code{get_supported_instruction_sets()}; \code{None} skips vectorization.
  \item[\code{assume_inner_stride_one}] replaces the innermost stride by 1, so that
    accesses along $x$ are contiguous. Kernels with runtime strides cannot be vectorized
    without it; the caller must guarantee it (waLBerla's fields in \code{fzyx} layout do).
  \item[\code{assume_aligned}] the first inner cell of each row is aligned, so some
    accesses may use aligned loads and stores (only for instruction sets that have them).
  \item[\code{nontemporal}] \code{True} or a set of fields whose stores should bypass
    the cache (\cref{r04:sec:nontemporal}).
  \item[\code{assume_sufficient_line_padding}] with aligned access: run the vector loop
    past the row end instead of adding a scalar remainder loop.
\end{description}
All fields of a kernel must have the same floating-point type; it decides between the
\code{float} and \code{double} variant of the instruction set.

\section{An instruction set is a dictionary}\label{r04:sec:dict}

\code{get_vector_instruction_set(data_type, name)}
(\file{backends/simd_instruction_sets.py}) returns a Python dictionary. Its keys name
operations, its values are C code with \code{str.format} placeholders. A few entries of
upstream's RVV dictionary for \code{float}:
\begin{center}\small
\begin{tabularx}{\linewidth}{@{}l>{\raggedright\arraybackslash}X@{}}
  \toprule
  \textbf{key} & \textbf{value}\\
  \midrule
  \code{'+'} & \code{__riscv_vfadd_vv_f32m1({0},{1},{loop_stop} - {loop_counter})}\\
  \code{'loadU'} & \code{__riscv_vle32_v_f32m1({0},{loop_stop} - {loop_counter})}\\
  \code{'makeVecConst'} & \code{__riscv_vfmv_v_f_f32m1({0}, {loop_stop} - {loop_counter})}\\
  \code{'width'} & \code{__riscv_vsetvlmax_e32m1()} (a C expression, not a number)\\
  \code{'float'} & \code{vfloat32m1_t}\\
  \code{'headers'} & \code{<riscv_vector.h>}, \code{"riscv_v_helpers.h"}\\
  \bottomrule
\end{tabularx}
\end{center}
The dispatcher picks the module by name: \code{neon}, \code{sme} and names starting with
\code{sve} go to the ARM module, \code{vsx} to PowerPC, \code{rvv} to the RISC-V module,
everything else to x86. Our \isa{rvv_ext} adds one branch there
(\cref{r04:sec:ext-rules}).

\section{The innermost loop}\label{r04:sec:loop}

\code{vectorize_inner_loops_and_adapt_load_stores} visits every innermost loop and
first decides what to do with the loop's end. The rule depends only on the dictionary:
\begin{enumerate}
  \item If the store template contains \code{{loop_stop}}, the loop is kept whole:
    the instruction set handles the last, partial vector itself. This is the case for
    the scalable instruction sets (\isa{rvv}, \isa{sve}), whose every intrinsic gets
    ``elements left'' as its active length.
  \item Otherwise, with \code{assume_aligned} and \code{assume_sufficient_line_padding},
    the loop end is rounded up to a multiple of the width.
  \item Otherwise the loop is cut at the last multiple of the width into a vector loop
    and a scalar remainder loop (\code{cut_loop}, which deep-copies the loop body). This
    needs an integer width; it is what SSE, AVX, NEON, \isa{sve256} and our
    \isa{rvv_ext-vls256} get.
\end{enumerate}
Then every field access whose index depends on the loop counter becomes a
\code{VectorMemoryAccess}: unit-stride, or strided if the instruction set has strided
loads and stores (RVV has). An access that is neither makes the vectorizer give up on
that loop with the warning ``Could not vectorize loop because of non-consecutive memory
access''; the loop stays scalar. Finally the loop step becomes the width, conditionals
become masks, and \code{insert_vector_casts} marks every scalar inside a vector
expression for broadcasting.

\section{Printing: templates filled in}\label{r04:sec:printer}

The C code is printed by \code{VectorizedCustomSympyPrinter}
(\file{backends/cbackend.py}), the same class for all instruction sets. For every
vector node it looks up the template and calls \code{format} on it with the already
printed operands:
\begin{pycode}
# _print_Add, simplified: positive terms first, then '+' or '-' templates
processed = summands[0].term
for summand in summands[1:]:
    func = self.instruction_set['-'] if summand.sign == -1 else self.instruction_set['+']
    processed = func.format(processed, summand.term, **self._kwargs)
\end{pycode}
\code{self._kwargs} holds \code{loop_counter} and \code{loop_stop} of the enclosing
loop; the loop printer sets them. That is how the RVV templates get their active length
\code{{loop_stop} - {loop_counter}} without the printer knowing anything about RVV.
A scalar marked for broadcasting is printed through \code{'makeVecConst'}, a number
through \code{'makeVecConst'} as well, an index vector through \code{'makeVecIndex'}.

\begin{keyidea}[Why this matters for our extras]
  The printer never looks inside a template; it only calls \code{format} on it. A
  template that is a subclass of \code{str} with its own \code{format} method can
  therefore decide, at print time, to print something else. Our \isa{vf} and
  \isa{fma} options are built on exactly this (\cref{r04:sec:ext-vf-fma}).
\end{keyidea}

\section{Nontemporal stores and cache-line zeroing}\label{r04:sec:nontemporal}

A store to a cache line that is not in the cache normally reads the line from memory
first (write-allocate). Instruction sets avoid that read in two ways, and the printer
supports both:
\begin{itemize}
  \item true nontemporal (streaming) store instructions, templates \code{'stream'}
    (x86, SVE) plus a fence after the kernel (\code{'streamFence'});
  \item zeroing the whole cache line before writing it, template \code{'cachelineZero'}
    (NEON, VSX and RVV), so that it never has to be read.
\end{itemize}
Zeroing is only safe if the whole line will be written and lies inside the array, so
for every nontemporal unit-stride store the printer emits a guard. This is the guard in
waLBerla-wind's generated kernel for one PDF direction:
\begin{ccode}
if (((uintptr_t) &_data_pdfs[_stride_pdfs_1*ctr_1 + _stride_pdfs_2*ctr_2 + ctr_0] & _clsize_mask) == 0
    && (_stride_pdfs_1*ctr_1 + _stride_pdfs_2*ctr_2 + ctr_0 + _clsize/4) < _size_pdfs_0*_size_pdfs_1*_size_pdfs_2) {
    cachelineZero((void*) &_data_pdfs[_stride_pdfs_1*ctr_1 + _stride_pdfs_2*ctr_2 + ctr_0]);
}
\end{ccode}
The first condition asks whether the address starts a cache line; the second whether a
whole line fits before the end of the array. \code{_clsize} and \code{_clsize_mask} are
computed once at the start of the kernel from the template \code{'cachelineSize'}. The
second condition multiplies the spatial sizes of the \emph{stored} field, so the
vectorizer sets \code{use_all_written_field_sizes} on the kernel, and every written
field's sizes become kernel parameters, even sizes the loops do not use. That flag is
the root of the waLBerla bug of \cref{r04:sec:blockwise-bug}.

\section{Detecting the hardware}\label{r04:sec:detect}

\code{get_supported_instruction_sets()} is what \code{'best'} and waLBerla's
\code{OPTIMIZE_FOR_LOCALHOST} rely on. On RISC-V Linux it reads \code{AT_HWCAP} with
\code{getauxval} and returns \code{['rvv']} if the bit for the V extension is set,
otherwise an empty list. The environment variable \code{PYSTENCILS_SIMD} (a
comma-separated list) overrides the detection on any machine. Detection never returns
\isa{rvv_ext}: our instruction set is used only where it is named.
